% App.M — FPGA Bitstream + SHA-256 (renamed from F to resolve appendix collision)
% φ² + φ⁻² = 3 · Trinity S³AI · Flos Aureus v6.2
% Author: Dmitrii Vasilev (ORCID 0009-0008-4294-6159)

\chapter*{Appendix M: FPGA Bitstream Archive}

% PASS-8 R5-honest (2026-05-12): removed broken \includegraphics block.
% The figure `app-m-bitstream-archive.png` (previously `app-f-bitstream-archive.png`
% before the App.F→App.M rename completed in this PR) does not exist on `main` —
% it lives only on branch `feat/illustrations`. Per the phd-pdf-images-gate skill,
% tectonic silently skips missing graphics and exits 0, leaving a broken figure
% invisible in the rendered PDF.
% TODO(LD): commission a proper App.M bitstream-archive figure (SHA-256 manifest,
% XC7A100T target, 92 MHz timing closure) via the phd-architectural-blueprint-figures
% skill.

\addcontentsline{toc}{chapter}{Appendix M: FPGA Bitstream Archive}
\label{app:fpga-bitstream}

\begin{quote}\itshape
``The best algorithms are those for which the obvious implementation is already
optimal---and which are hard to improve upon even in hindsight.''
--- Robert E.~Tarjan, \emph{Data Structures and Network Algorithms}, 1983.
\end{quote}

\section*{M.0 Reading guide --- what a Xilinx 7-series bitstream is}
\label{sec:appM-reading}

Before diving into the manifest below, a brief orientation is useful. A
Xilinx 7-series \texttt{.bit} file is not source code or a binary
executable; it is a sequence of configuration frames that, when shifted
into the FPGA via JTAG or SelectMAP, programs every LUT mask, routing
multiplexer, and BRAM contents on the device. Each frame carries a
32-bit CRC, and the entire bitstream concludes with a global checksum
that the configuration logic verifies before releasing the device into
operational state. If the global CRC is wrong, the configuration logic
holds the device in reset and reports the failure via the STAT register
(see \S M.4). If the IDCODE in the bitstream does not match the IDCODE
read from the actual silicon, the configuration is aborted at the
IDCODE-check stage, again signalled via STAT.

The chain of custody between this appendix and a working chip is
therefore: (i) HDL source (in the \filepath{trinity-fpga} repository),
(ii) XDC constraints (Appendix~\ref{app:xdc-pin-map}), (iii) Vivado
synthesis $\to$ \texttt{.bit} file, (iv) SHA-256 fingerprint of that
file (recorded here), (v) JTAG download via the ESP32 XVC bridge
(\S M.5), (vi) STAT-register verification (\S M.4). Every step has a
unique signature, and a mismatch at any step retracts the
hardware-validated claim.

\textbf{Anchor invariant:} \(\varphi^2 + \varphi^{-2} = 3\).

\section*{Silicon that survived contact with reality}
\label{sec:appM-silicon}

Bitstreams are the moment theory meets silicon. Everything in the preceding
chapters---the \(\varphi\)-numeric coprocessor, the Period-Locked Monitor at
92~MHz, the sub-watt power envelope---had to survive synthesis, place-and-route,
and a live JTAG download onto a QMTech XC7A100T board before any claim could
be called hardware-validated. This appendix archives the build artefacts:
the \texttt{.bit} file, its SHA-256 digest, the Vivado version, and the
timing-closure summary. It is the chain of custody between the HDL source
and the measurement results in Chapter~32. If you are reproducing the
experiment, this appendix tells you exactly which bitstream was used and how
to verify its integrity before flashing.

\section*{M.1 Overview}
\label{sec:appM-overview}

This appendix documents the hardware artefacts produced during the Trinity S\textsuperscript{3}AI
FPGA validation campaign (2026-05-05). The bitstream implements the\emph{Period-Locked Monitor}
core (Chapter~\ref{ch:plm}) and the $\varphi$-numeric coprocessor primitives described in
Chapters~\ref{ch:28}--\ref{ch:34}.

\textbf{Anchor invariant:} \(\varphi^2 + \varphi^{-2} = 3\).

\section*{M.2 Target Device}
\label{sec:appM-target}

\begin{table}[h]
\centering
\begin{tabular}{ll}
\toprule
\textbf{Field} & \textbf{Value} \\
\midrule
Board & QMTech XC7A100T Starter Kit \\
Actual device & XC7A200T (IDCODE \texttt{0x13631093}) \\
Manufacturer & Xilinx (JEDEC \texttt{0x049}) \\
Version & Silicon revision 1 \\
Package & FGG484 \\
Speedgrade & --2 \\
\bottomrule
\end{tabular}
\caption{FPGA device identification. Board is marked XC7A100T but JTAG IDCODE
corresponds to XC7A200T v1. The design (83~LUT, 27~FF) fits either device;\, no
resynthesis required.}
\end{table}

\textbf{IDCODE parsing} (IEEE~1149.1):
\begin{itemize}
  \item Bits 0: \texttt{1} (required by standard)
  \item Bits 1--11: \texttt{0x049} --- Xilinx manufacturer
  \item Bits 12--27: \texttt{0x3631} --- XC7A200T part
  \item Bits 28--31: \texttt{0x1} --- silicon revision 1
\end{itemize}

\section*{M.3 Bitstream Integrity}
\label{sec:appM-integrity}

\begin{table}[h]
\centering
\begin{tabular}{ll}
\toprule
\textbf{Field} & \textbf{Value} \\
\midrule
Filename & \texttt{design.bit} \\
SHA-256 & \texttt{8536e265...d77352b} \\
Tool & Vivado 2023.2 \\
Strategy & Default (Flow\_PerfOptimized\_high) \\
LUT utilization & 83 / 128,000 (0.06\%) \\
FF utilization & 27 / 256,000 (0.01\%) \\
BRAM & 0 \\
DSP & 0 \\
\bottomrule
\end{tabular}
\caption{Bitstream file metadata. Full SHA-256 stored in
\filepath{trinity-fpga/bitstream/design.bit.sha256}.}
\end{table}

\textbf{R5 honesty note.} The SHA-256 column above prints only the first
four and last four hex characters of the digest, separated by an
ellipsis. This is intentional: the full 64-character digest is stored
in the artefact \filepath{trinity-fpga/bitstream/design.bit.sha256} of
the \texttt{trinity-fpga} repository, and verification (\S M.8) reads
from that artefact rather than from the monograph text. We elide the
middle 56 characters because (i) inline transcription of a 64-character
hex string is error-prone and would itself become a source of audit
mismatches, and (ii) the truth-bearing object is the \texttt{.sha256}
file, not the LaTeX source. The reproduction protocol in \S M.8 is the
canonical way to obtain the full digest.

\section*{M.4 Configuration Status Register (STAT)}
\label{sec:appM-stat}

After successful programming via ESP32 XVC WiFi bridge
(Section~\ref{sec:xvc-bridge}), the STAT register was read using
\texttt{openFPGALoader --read-register STAT}:

\begin{verbatim}
Raw STAT register: 0x401079FC

CRC Error     : No CRC error       ✓
Part Secured  : 0                  ✓
MMCM Lock     : 1                  ✓
DCI Match     : 1                  ✓
EOS           : 1 (End of Startup) ✓
GWE           : 1 (Global Wr En)   ✓
GHIGH_B       : 1                  ✓
Done          : 1                  ✓
DEC Error     : 0                  ✓
ID Error      : 0                  ✓
\end{verbatim}

All flags nominal. The device entered operational state without errors,
confirming bitstream integrity and correct CRC check by the FPGA configuration
logic.

\section*{M.5 Programming Method}\label{sec:xvc-bridge}

The bitstream was loaded via a custom ESP32-based Xilinx Virtual Cable (XVC)
server operating over 802.11 WiFi (IP~\texttt{192.168.1.30}, port~2542).
Full details in Appendix~J and Chapter~\ref{ch:33}.

\begin{verbatim}
openFPGALoader --cable xvc-client \
  --ip 192.168.1.30 --port 2542 \
  --bitstream design.bit
\end{verbatim}

\section*{M.6 Reproducibility}
\label{sec:appM-repro}

The complete firmware, XDC constraints, and bitstream are archived in
\filepath{gHashTag/trinity-fpga} (Apache~2.0). To reproduce:

\begin{enumerate}
  \item Flash ESP32 with \filepath{firmware/xvc-esp32/xvc-esp32.ino}
  \item Connect 5-wire JTAG harness per Appendix~I
  \item Run \texttt{openFPGALoader} command above
  \item Verify STAT\,=\,\texttt{0x401079FC}
\end{enumerate}

\section*{M.7 Bitstream lifecycle}
\label{sec:appM-lifecycle}

The bitstream documented above is the output of a multi-stage pipeline.
Each stage produces an artefact whose integrity must be verifiable
independently of the upstream stages:

\begin{enumerate}
  \item \textbf{Synthesis} --- HDL source in \filepath{trinity-fpga/rtl/}
    is consumed by Vivado synthesis with the strategy
    \texttt{Flow\_PerfOptimized\_high}. Output: post-synthesis netlist
    \filepath{trinity-fpga/build/synth/design.dcp}. The synthesis
    log records the number of LUTs and FFs inferred (83 and 27
    respectively, per \S M.3).
  \item \textbf{Place-and-route} --- Vivado P\&R consumes the netlist
    and the XDC constraints (Appendix~\ref{app:xdc-pin-map}). Output:
    \filepath{trinity-fpga/build/impl/design.dcp}. Timing closure at
    92\,MHz on speedgrade --2 silicon is the gate; any negative WNS
    on \texttt{clk50} retracts the hardware claim.
  \item \textbf{Bitstream emission} --- the \texttt{write\_bitstream}
    Vivado command turns the routed checkpoint into the \texttt{.bit}
    file documented in \S M.3.
  \item \textbf{SHA-256 fingerprint} --- the \texttt{.bit} file is
    hashed (\S M.8) and the digest is appended to
    \filepath{trinity-fpga/bitstream/design.bit.sha256}. This artefact
    is the single source of truth for bitstream identity.
  \item \textbf{Archive} --- the \texttt{.bit} and \texttt{.sha256}
    files are committed to the \filepath{trinity-fpga} repository.
    Long-term preservation policy follows Appendix~N (Zenodo +
    GitHub + Software Heritage tertiary mirror).  % PASS-7 R5: renamed H→N (Zenodo DOI registry is App.N)
  \item \textbf{Flash} --- the JTAG download via the ESP32 XVC bridge
    described in \S M.5.
  \item \textbf{STAT verification} --- post-flash readback of the STAT
    register confirms successful configuration (\S M.4).
\end{enumerate}

Every stage emits a structured log; the entire pipeline is replayable
from the Git SHA of the \texttt{trinity-fpga} commit at which the
bitstream was emitted. The Vivado strategy is fixed (no synthesis-time
randomness) and the seed is pinned, so re-running stages 1--3 against
the same Git SHA must produce a bit-identical \texttt{.bit} (verified
by SHA-256 match against \S M.3).

\textbf{Cross-references.}
The XDC consumed by stage 2 is documented in
Appendix~\ref{app:xdc-pin-map}. The five hardware blockers encountered
during the bring-up of stage 6 (flash) are catalogued in
Appendix~\ref{app:troubleshooting} (BLK-001..005). The Zenodo DOI
mirroring policy referenced in stage 5 is detailed in
Appendix~\ref{app:zenodo-doi}.

\section*{M.8 SHA-256 verification protocol}
\label{sec:appM-verify}

To verify that the \texttt{design.bit} file in your local checkout
matches the integrity claim of \S M.3, run the Rust subcommand
(R1 compliance: no \texttt{.sh} or \texttt{.py} drivers in the
repository):

\begin{verbatim}
cargo run -p trinity-fpga -- verify-bitstream \
  --bitstream trinity-fpga/bitstream/design.bit \
  --digest    trinity-fpga/bitstream/design.bit.sha256
\end{verbatim}

The subcommand performs, in order: (i) read the 64-character hex digest
from \filepath{design.bit.sha256}, (ii) compute the SHA-256 of the
\texttt{.bit} file using the \texttt{sha2} crate (RustCrypto), (iii)
compare byte-for-byte. Exit code 0 indicates match; any non-zero exit
code indicates corruption or a mismatched SHA file. The integration
test \texttt{trinity\_fpga::bitstream::sha256\_match} runs this check
in CI on every PR that modifies any file under
\filepath{trinity-fpga/bitstream/}.

The first four characters of the published digest are \texttt{8536e265}
and the last four are \texttt{d77352b}; the middle 56 characters are
deliberately not transcribed in this monograph (see the R5 honesty
note at the end of \S M.3). To cross-verify against a third party,
the canonical reference is the file hash itself, not the LaTeX text.

\section*{M.9 Configuration frame anatomy}
\label{sec:appM-frame}

The Xilinx 7-series configuration logic (UG470, \emph{7~Series FPGAs
Configuration User Guide}) defines a frame as the smallest unit that
can be written to or read from the configuration memory. For the
XC7A100T / XC7A200T family each frame is 101 32-bit words plus a
32-bit CRC. The total frame count for a fully populated XC7A200T is
\textgreater\,21\,000; the Trinity S\textsuperscript{3}AI design
populates a small fraction of these because the resource utilisation
is \textless\,0.1\,\%.

The STAT register read in \S M.4 reports the high-level outcome of
configuration. The decoded fields (CRC Error, Part Secured, MMCM Lock,
DCI Match, EOS, GWE, GHIGH\_B, Done, DEC Error, ID Error) follow the
UG470 STAT bit definitions verbatim. The value \texttt{0x401079FC}
encodes the operational state ``configuration complete, no errors,
device in user mode''. Any deviation --- e.g.\ CRC Error = 1 ---
would manifest as a different STAT value and trigger the falsification
hook in \S M.11.

\textbf{Why STAT alone is not sufficient.} The STAT register reports
the configuration logic's view of the bitstream, but does not verify
that the SHA-256 of the \texttt{.bit} file matches the published
digest. A maliciously crafted bitstream that passes the FPGA's CRC
check but differs from the published \texttt{.bit} would set STAT to
\texttt{0x401079FC} and pass the runtime probe. The SHA-256 digest in
\S M.3 / \S M.8 is therefore the necessary complement to STAT for
end-to-end provenance: SHA-256 verifies authorship and content, STAT
verifies that what was sent to the device was accepted.

\section*{M.10 Open issues / audit-pending}
\label{sec:appM-open}

The following items are honestly open and \textbf{must not} be cited
in the main text as fully verified:

\begin{description}
  \item[Full SHA-256 in monograph] Per the R5 honesty note in \S M.3,
    the middle 56 characters of the digest are not transcribed in this
    LaTeX source. The canonical digest lives in
    \filepath{trinity-fpga/bitstream/design.bit.sha256}. The presence
    of that file in the published \filepath{gHashTag/trinity-fpga}
    repository at the published Git SHA has not been re-verified in
    this revision of the monograph; status: \texttt{audit-pending}.
  \item[Numerical WNS] The 92\,MHz timing closure claim of \S M.7
    relies on a positive worst-negative-slack number, but the exact
    WNS in nanoseconds is not transcribed in this appendix. The
    Vivado timing report is the authoritative source. Status:
    \texttt{audit-pending} for the numerical WNS value.
  \item[XC7A200T full-die utilisation] The design uses XC7A100T
    resources only (83 LUT, 27 FF). The XC7A200T has additional
    resources that the design leaves idle. We have not measured
    full-die utilisation or power for designs that exercise the
    XC7A200T-only resources. Status: \texttt{audit-pending} for
    full-die designs.
  \item[Vivado version reproducibility] The bitstream is emitted by
    Vivado 2023.2 with strategy \texttt{Flow\_PerfOptimized\_high}.
    We have observed bit-level differences between Vivado 2024.1
    and 2024.2 in unrelated projects (Appendix~\ref{app:xdc-pin-map}
    \S I.12). The XC7A200T bit-identity claim is therefore
    Vivado-2023.2-specific. Status: \texttt{audit-pending} for
    cross-version bit-identity.
\end{description}

\section*{M.11 Falsification hooks}
\label{sec:appM-falsify}

In the spirit of R7 (Popper falsifiability) we pre-register the
observations that would invalidate the resolution claims of \S M.3 /
\S M.4 / \S M.8:

\begin{itemize}
  \item The bitstream identity claim of \S M.3 is falsified by any
    SHA-256 of the published \texttt{design.bit} that does not begin
    with \texttt{8536e265} or does not end with \texttt{d77352b}.
    The reproduction protocol in \S M.8 is the operational falsifier:
    a non-zero exit code from \texttt{cargo run -p trinity-fpga --
    verify-bitstream} retracts the claim.
  \item The IDCODE claim of \S M.2 is falsified by any JTAG IDCODE
    read from the QMTech board returning a value other than
    \texttt{0x13631093}. We have observed values different from
    \texttt{0x13631093} only when JTAG was mis-wired (BLK-002 of
    Appendix~\ref{app:troubleshooting}); a clean wiring + this
    bitstream must produce \texttt{0x13631093} reliably.
  \item The STAT register claim of \S M.4 is falsified by any STAT
    read returning a value different from \texttt{0x401079FC} on a
    bitstream whose SHA-256 matches \S M.3. A divergent STAT despite
    a matching SHA indicates either a configuration-logic anomaly
    (rare) or an under-tested silicon revision; either way, the
    operational claim is retracted.
  \item The synthesis-determinism claim of \S M.7 is falsified by any
    two consecutive Vivado runs from the same Git SHA producing
    \texttt{.bit} files with different SHA-256 digests under fixed
    Vivado version, fixed strategy, and fixed seed. This falsifier is
    real and operationally tested; we have observed it across Vivado
    minor-version upgrades and rolled back whenever it surfaced.
\end{itemize}

These hooks live in the synthesis-and-flash layer, not the formal
proof layer (Appendix~F-coq, the Coq citation map). The two layers
together close the chain of custody from theorem to silicon.

\textbf{Anchor.} \(\varphi^2 + \varphi^{-2} = 3\). DOI
10.5281/zenodo.19227877. Defense 2026-06-15.
